This section develops the protocol in full: the actors and threat model, the canary
construction that manufactures the null, the filtration that makes the martingale
argument literally correct, the certified quantity, the betting engine with its
explicit wealth dynamics, and finally the assembled protocol as pseudocode
(Algorithm~\ref{alg:vouch}) and as a flowchart (Figure~\ref{fig:flowchart}).

\subsection{Actors, pipeline, and threat model}\label{sec:setting}

The provider fine-tunes a base model on
$D_{\text{keep}} \cup D_{\text{forget}} \cup C$, where $C$ is the canary set described
below. When deletion requests arrive, an unlearning algorithm $\Ucal$---negative
preference optimisation, GradDiff, and so on---produces an unlearned model $M_u$ from
the fine-tuned model, using the forget set augmented with the relevant canaries. A
verifier, holding only black-box query access to $M_u$ together with a canary manifest
whose commitment was published before unlearning began, runs VOUCH and either issues,
withholds, or revokes a certificate.

Throughout we assume the \emph{honest-but-verifiable} setting: the provider genuinely
executes $\Ucal$ on the forget set as declared, and the open question is whether it
worked and how well, under quantified error control. This assumption is doing real
work and we state it wherever the certification claim is summarised, because the
commitment prevents the manifest from being changed after the fact but does not
prevent a provider who has \emph{located} the cohort from treating it differently
from organic data. Section~\ref{sec:exp-detect} measures how hard locating the cohort
actually is---the answer is: much easier than we had assumed---and explains why
locating it is nonetheless not the same as defeating the audit: both twins in a pair
are equally conspicuous, so a provider who finds the cohort still learns nothing
about which twin the coin admitted. A provider willing to special-case the committed
canaries is outside our scope and, as prior work establishes, can only be addressed
by cryptographic attestation \citep{eisenhofer2025verifiable}, with which VOUCH
composes.

\subsection{Paired ghost canaries}\label{sec:canaries}

A templated generator emits exchangeable twin pairs
$(c_i^0, c_i^1)$ for $i = 1,\dots,m$. The two twins in a pair are independent draws from
the same conditional template distribution---the same surface form populated with
independently sampled fictitious entities carrying secrets of at least $40$ bits of
entropy---so that within a pair they are exchangeable by construction. We match the
template to the host corpus, using question--answer pairs for TOFU and record-style
sentences for news text. For each pair an
independent fair coin $b_i \sim \Bern(\tfrac12)$ selects the \emph{in-twin}
$c_i^{b_i}$, which we insert into
the fine-tuning corpus with a repetition factor $r_i \in \{1,2,4,8\}$ (these strata give
us a dose--response curve) and later append to the forget set, so that it traverses the
unlearning pipeline exactly as organic forget data does. The \emph{ghost twin}
$c_i^{1-b_i}$ is never shown to any model.

\subsection{Commitment and filtration}\label{sec:filtration}

Before unlearning begins the provider publishes a commitment to the manifest. The
commitment is a Merkle tree whose $i$-th leaf is
\begin{equation}\label{eq:commit}
\ell_i \;=\; H\bigl(c_i^0,\, c_i^1,\, b_i,\, r_i,\, \rho_i\bigr),
\qquad H = \text{SHA-256},
\end{equation}
with $\rho_i$ a per-leaf blinding nonce derived from the manifest seed, and the
published value is the tree root. This binds the pairs and the coins and prevents the
audit from being gamed after the fact, exactly as a flat hash of the whole manifest
would; the tree structure buys one further thing that a flat hash does not, and the
theory needs it. At step $t$ the provider hands over leaf $\pi(t)$ together with its
authentication path, which the verifier checks against the published root; the path is
$\lceil\log_2 m\rceil$ hashes long, so per-pair opening costs nine hashes at
$m = 384$. The blinding nonce matters because the template library is public: without
it, an unopened leaf's hash could be attacked by enumeration.

The verifier draws a uniformly random reveal order $\pi$, committed in advance, and
opens leaf $\pi(t)$ at step $t$, checking its Merkle path against the published root.
The filtration is
\begin{equation}\label{eq:filtration}
\F_t \;=\; \sigma\Bigl(\text{query access to } M_u,\;
\bigl\{(c_{\pi(u)}^0, c_{\pi(u)}^1, b_{\pi(u)}, r_{\pi(u)})\bigr\}_{u \le t},\;
\text{tie coins } 1..t \Bigr).
\end{equation}
Opening leaf by leaf rather than all at once is what makes $b_{\pi(t)}$ genuinely
unrevealed at time $t-1$, and hence what makes ``the coin is still fair given the
past'' a true statement about $\F_{t-1}$ rather than an informal one. Every
supermartingale argument in Section~\ref{sec:theory} is with respect to
\eqref{eq:filtration}. Nothing else about the protocol changes, and the verifier's
computational cost is unaffected.

\subsection{Scores and the certified quantity}\label{sec:scores}

We fix a declared class $\F$ of scalar score functions, each aggregated over $Q$ fixed
query prompts, and each oriented so that a \emph{higher} score means \emph{more
memorised}:
\begin{itemize}
\item $s_{\mathrm{loss}}$, the token-normalised \emph{log-likelihood} of the secret
      span---that is, the negative of its token-normalised negative log-likelihood;
\item $s_{\mathrm{mink}}$, the mean of the lowest $k\%$ token log-probabilities of the
      span ($k = 20$);
\item $s_{\mathrm{ratio}}$, the base-model-calibrated difference
      $s_{\mathrm{loss}}(M_u,\cdot) - s_{\mathrm{loss}}(M_0,\cdot)$, available when the
      verifier may query the pristine base model;
\item optionally $s_{\mathrm{probe}}$, a linear probe on the hidden state, calibrated
      on a disjoint cohort.
\end{itemize}
The sign convention matters and we fix it once, since an earlier draft described
$s_{\mathrm{loss}}$ as a negative log-likelihood while the released implementation
uses the log-likelihood, which flips the sign of every reported gap. For pair $i$ and
score $s$ we form the within-pair difference and its sign,
\begin{equation}\label{eq:sign}
D_i^{(s)} = s(M_u, c_i^{\text{in}}) - s(M_u, c_i^{\text{ghost}}), \qquad
Z_i^{(s)} = \mathbf{1}\{D_i^{(s)} > 0\},
\end{equation}
breaking ties by a committed fair coin. With scores oriented as above, a
\emph{positive} $D$ means the trained twin is the better-remembered one, which is the
direction of leakage; the worked example of Figure~\ref{fig:example} reports
$D = +8.0$ for a fine-tuned model whose in-secret costs $0.99$ nats per token against
the ghost's $9.03$, and this is now consistent with \eqref{eq:sign}, with the figure,
with the tables, and with the code. Exact ties are rare but not always negligible: we
measure a rate of $0$ in $124{,}416$ scored pairs for $s_{\mathrm{loss}}$ across every
architecture, and $2.7\%$ for $s_{\mathrm{mink}}$ on Gemma-4-E2B, where the min-$k\%$
statistic over a short span takes few distinct values
(Section~\ref{sec:exp-bench}). The committed tie-break coin is what keeps the null
exact in that cell.

The quantity we certify is $\sup_{s\in\F}|\Delta^{(s)}|$ with $\Delta^{(s)}$ the
realised cohort advantage of Definition~\ref{def:delta}. We are deliberate that this
is influence extractable \emph{through $\F$} on the \emph{committed canary cohort}: it
is the strongest target compatible with the impossibility results for behavioural
audits, and we make no claim about erasure in parameter space, about scores outside
$\F$, or---without the explicit inflation of Proposition~\ref{prop:transfer}---about
the super-population advantage.

\subsection{The betting engine}\label{sec:engine}

All three of VOUCH's outputs are built from one primitive, the \emph{betting
e-process}, and it is worth writing its dynamics out explicitly because the validity
proofs in Section~\ref{sec:theory} reduce to inspecting them. A gambler starts with
unit wealth and, before the pair opened at step $t$ is revealed, stakes a predictable
fraction $\lambda_t$ (chosen from the information $\F_{t-1}$ revealed so far) on the
sign $Z_{\pi(t)}$ deviating from a hypothesised mean.

For the \emph{certificate} arm the hypothesised mean is not a constant but the
predictable without-replacement boundary $p_0(t)$ of Eq.~\eqref{eq:worboundary}: the
mean the unrevealed remainder of the cohort would have to carry for the null
$\Delta^{(s)} \ge \eps$ still to hold. Against that boundary and its mirror the wealth
processes for score $s$ are
\begin{equation}\label{eq:wealthwor}
E_t^{s,+} = \prod_{u \le t}\bigl(1 + \lambda_u^{s,+}\,(p_0(u) - Z^{(s)}_{\pi(u)})\bigr),
\qquad
E_t^{s,-} = \prod_{u \le t}\bigl(1 + \lambda_u^{s,-}\,(Z^{(s)}_{\pi(u)} - \bar p_0(u))\bigr),
\end{equation}
with $\bar p_0(t)$ the analogous boundary for the null $\Delta^{(s)} \le -\eps$ and
stakes confined to the admissible range so that wealth can never go negative. Because
sequential reveal of a uniform permutation is simple random sampling without
replacement, the conditional mean of the next sign is exactly the mean of the
remainder, so under the null each factor has conditional mean at most one: each
process is a nonnegative supermartingale started at one, and Ville's inequality
converts wealth into evidence. Wealth above $1/\alpha$ is a licence to reject at
whatever time we happen to be looking. Setting $p_0(t) \equiv p_0$ instead recovers
the fixed-boundary process of the earlier draft, which targets the super-population
advantage and requires the conditional-mean null; we retain it as a reported variant
and compare the two throughout Section~\ref{sec:exp}.

For the \emph{revocation} arm the hypothesised mean is the constant $\tfrac12$, which
Theorem~\ref{thm:null} makes exact conditionally, so no boundary recentring is needed.

How the stakes are chosen determines power, not validity, which frees us to be
aggressive. VOUCH plays a uniform mixture of eight experts and lets their wealth
arbitrate: six constant-fraction bettors $\lambda = c \cdot \lambda_{\max}$ with
$c \in \{0.02, 0.05, 0.1, 0.2, 0.4, 0.7\}$, an online-Newton-step (ONS) learner, and a
truncated Krichevsky--Trofimov (KT) plug-in. The ONS expert follows the coin-betting
recursion
\begin{equation}\label{eq:ons}
g_t = \frac{p_0(t) - Z_t}{1 + \lambda_t (p_0(t) - Z_t)}, \qquad
A_t = A_{t-1} + g_t^2, \qquad
\lambda_{t+1} = \Pi_{[0,\lambda_{\max}]}\!\Bigl(\lambda_t + \tfrac{\eta\, g_t}{A_t}\Bigr),
\end{equation}
with $A_0 = 1$, $\eta = \tfrac12$, and $\Pi$ the clip onto the admissible range; its
log-wealth regret against the best constant bet in hindsight is $O(\log t)$. The KT
expert bypasses stake selection altogether: it estimates the mean by
$q_t = (\sum_{u<t} Z_u + \tfrac12)/t$, truncates $q_t$ to the alternative side of
$p_0(t)$, and multiplies wealth by the plug-in likelihood ratio
\begin{equation}\label{eq:kt}
e_t \;=\; Z_t\,\frac{q_t}{p_0(t)} \;+\; (1 - Z_t)\,\frac{1 - q_t}{1 - p_0(t)},
\end{equation}
which has conditional mean at most one under the null and achieves the optimal growth
rate up to a $\tfrac1{2t}\log t$ redundancy. Because a convex combination of
e-processes with a fixed prior is again an e-process, the mixture inherits exact
validity while paying only a $\log 8$ additive cost in log-wealth relative to its best
expert---power adapts online with nothing to tune.

The same wealth idea, run over a grid rather than at a single boundary, yields the
quantitative face of the certificate. For every candidate mean $m'$ on a grid of
$1001$ points in $[0,1]$ we maintain two wealth processes betting that the truth is
above and below $m'$, each recentred by the same without-replacement rule, with
empirical-Bernstein (aGRAPA) stakes
\begin{equation}\label{eq:agrapa}
\lambda_t(m') \;=\; \Pi_{[0,\,\lambda_{\max}]}
\Bigl(\tfrac{\hat\mu_{t-1} - m'}{\hat\sigma_{t-1}^2 + (\hat\mu_{t-1} - m')^2}\Bigr),
\end{equation}
where $\hat\mu_{t-1}$ and $\hat\sigma^2_{t-1}$ are the running mean and variance of the
signs. The set of grid values whose wealth on both sides stays below $2/\alpha$ and
whose recentred boundary remains in $[0,1]$, intersected over time,
\begin{equation}\label{eq:cs}
C_t \;=\; \bigcap_{u \le t}\,\bigl\{\, m' :\ \max\bigl(W_u^{\uparrow}(m'),\,
W_u^{\downarrow}(m')\bigr) < 2/\alpha \,\bigr\},
\end{equation}
is a $(1-\alpha)$ confidence sequence for the cohort sign rate: it covers the truth at
every time simultaneously, so the auditor may glance at it whenever they please, and
it collapses onto the exact cohort value once the cohort is exhausted. Mapped through
$\Delta = 2p - 1$ it reports the residual advantage with its anytime-valid
uncertainty; Remark~\ref{rem:cs} states precisely which coverage claims are per-score
and which are simultaneous over $\F$.

The revocation arm inverts the role of the null. Under \emph{exact} unlearning,
Theorem~\ref{thm:null} makes every $Z_{\pi(t)}^{(s)}$ an exact conditional fair coin,
so here---and only here---adaptive score weighting is sound. We run, per score and per
direction, e-processes against $\Efrak[Z] = \tfrac12$ and combine them with predictable
wealth-proportional weights: with $w_{s,t} \propto \exp(\log E^{s}_{t-1})$ computed
before the pair at step $t$ is opened,
\begin{equation}\label{eq:revmix}
R_t \;=\; \prod_{u\le t}\;\sum_{s\in\F} w_{s,u}\,
\frac{E^{s}_{u}}{E^{s}_{u-1}},
\end{equation}
which is an e-process because each factor is a predictable convex combination of
conditional-mean-one increments. When $R_t \ge 1/\alpha$ the verifier holds
anytime-valid evidence of residual memorisation and the certificate is withheld or
revoked; the weights meanwhile steer wealth toward whichever score discriminates best,
which is exactly what an alarm should do.

\subsection{The protocol end to end}\label{sec:protocol}

Algorithm~\ref{alg:vouch} assembles the pieces into the three-phase protocol, and
Figure~\ref{fig:flowchart} shows the same pipeline as a flowchart. Phase~0 runs at
design time: generate the cohort, flip the coins, publish the Merkle root, fine-tune
with the in-twins planted. Phase~1 is the provider's unlearning step, in which the
in-twins ride along with the organic forget set. Phase~2 is the audit: stream the
pairs in the committed random order, opening and checking one leaf at a time, updating
the certificate processes, the confidence sequences, and the revocation mixture after
each pair, and acting on whichever boundary is crossed first. R-VOUCH then re-runs the
same loop on probed variants of $M_u$---after relearning on public data, after 4-bit
quantisation, and under adversarial prompt wrappers---and the robust certificate
requires every probe to pass. Our empirical probe coverage is partial and we state it
here rather than leave it to be inferred: P1 and P3 run on every benchmark cell and
every seed but are anchored on NPO; P2 runs on the GPT-2 and TinyGPT tiers only,
because the shared-frozen-base multi-adapter design would have to materialise merged
weights to quantise a multi-billion-parameter model. Streaming deletion waves each draw a fresh cohort with
independent coins; Theorem~\ref{thm:streaming} dictates how their verdicts may be
combined.
